\chapter{User Interface Design}

\section{Two Dashboard Applications}

The project uses two dashboard applications because internal operators and
merchant users have different permissions, risks, and workflows.

\begin{table}[H]
\centering
\caption{Dashboard split}
\begin{tabularx}{\textwidth}{L{0.24\textwidth}L{0.24\textwidth}Y}
\toprule
\textbf{UI} & \textbf{Audience} & \textbf{Purpose} \\
\midrule
Ops Dashboard & Internal Admin/Ops & Operate merchants, investigate evidence, manage users, and recover failed workflows. \\
Merchant Dashboard & Merchant dashboard users & Read merchant-scoped payment, refund, webhook, profile, credential, and analytics data. \\
\bottomrule
\end{tabularx}
\end{table}

This split is not only visual. It mirrors the access model of the system:
internal session cookies and RBAC for Ops, and merchant-scoped portal sessions
for dashboard users. The result is a clearer user experience and stronger data
separation.

\section{Ops Dashboard UI}

The Ops Dashboard is designed for repeated operational work. It favors dense but
readable list and detail layouts, explicit lifecycle actions, semantic badges,
and direct access to operational evidence such as onboarding data, credentials,
payments, refunds, webhook attempts, audit logs, and reconciliation records.

\begin{figure}[H]
\centering
\safeimage{ops-overview.png}{Ops Dashboard overview showing internal metrics and work queues.}
\caption{Ops Dashboard overview. The landing page is optimized for operational
prioritization and queue awareness.}
\end{figure}

\begin{figure}[H]
\centering
\safeimage{ops-merchant-detail.png}{Ops merchant detail with lifecycle status, credential metadata, portal users, onboarding evidence, recent payments, and audit.}
\caption{Ops merchant detail. Merchant status, onboarding evidence, credential
metadata, and portal-user provisioning are grouped into one internal workflow.}
\end{figure}

\section{Merchant Dashboard UI}

The Merchant Dashboard is read-only except for local password change. It gives
merchant users an overview first, then deeper analytics and explorer pages for
support or business questions. Long values such as webhook URLs, QR payloads,
and JSON snippets are wrapped or scrollable to avoid layout overflow.

\begin{figure}[H]
\centering
\begin{subfigure}{0.48\textwidth}
  \centering
  \safeimage[\linewidth]{merchant-overview.png}{Merchant Overview with summary cards and compact charts.}
  \caption{Overview}
\end{subfigure}
\hfill
\begin{subfigure}{0.48\textwidth}
  \centering
  \safeimage[\linewidth]{merchant-analytics.png}{Merchant Analytics with interactive Recharts, tooltips, and data tables.}
  \caption{Analytics}
\end{subfigure}
\caption{Merchant Dashboard overview and analytics.}
\end{figure}

\begin{figure}[H]
\centering
\begin{subfigure}{0.48\textwidth}
  \centering
  \safeimage[\linewidth]{merchant-payments.png}{Merchant payment explorer with filters and detail panel.}
  \caption{Payments}
\end{subfigure}
\hfill
\begin{subfigure}{0.48\textwidth}
  \centering
  \safeimage[\linewidth]{merchant-webhooks.png}{Merchant webhook explorer with event status and attempts.}
  \caption{Webhooks}
\end{subfigure}
\caption{Merchant explorer pages.}
\end{figure}

\begin{figure}[H]
\centering
\begin{subfigure}{0.48\textwidth}
  \centering
  \safeimage[\linewidth]{merchant-profile.png}{Read-only merchant profile and integration metadata.}
  \caption{Profile}
\end{subfigure}
\hfill
\begin{subfigure}{0.48\textwidth}
  \centering
  \safeimage[\linewidth]{merchant-credentials.png}{Credential metadata without raw secret.}
  \caption{Credentials}
\end{subfigure}
\caption{Read-only merchant metadata pages.}
\end{figure}

\section{Responsive And State Design}

Both dashboards treat asynchronous state as a first-class design concern:

\begin{itemize}
  \item loading, error, and empty states are explicit on data-fetching routes;
  \item long tables and side panels remain scrollable instead of forcing layout
  overflow;
  \item chart data is paired with readable labels and tables so exact values are
  not hidden behind hover-only interactions;
  \item session and status panels remain visually stable to reduce user
  confusion during multi-step workflows;
  \item the merchant login form uses explicit labels and accessible field
  associations.
\end{itemize}
